%======================================================================
% فصل اول: مقدمه
%======================================================================
\chapter{مقدمه}
\label{ch:intro}

%----------------------------------------------------------------------
\section{کلیات}
\label{sec:intro-overview}
%----------------------------------------------------------------------

در دهه‌ی اخیر، \term{هوش مصنوعی}{Artificial Intelligence (AI)} تحولی
بنیادین در حوزه‌های مختلف فناوری اطلاعات ایجاد کرده است. ظهور
\term{مدل‌های زبانی بزرگ}{Large Language Models (LLM)} مانند
\lr{GPT-4o}، \lr{Gemini} و \lr{Claude}، امکانات بی‌سابقه‌ای را برای
توسعه‌دهندگان نرم‌افزار و سازمان‌ها فراهم کرده است. این مدل‌ها از طریق
\term{رابط‌های برنامه‌نویسی کاربردی}{Application Programming Interface (API)}
در دسترس عموم قرار گرفته‌اند و کاربردهای گسترده‌ای در حوزه‌هایی از جمله
دستیاران هوشمند، تولید محتوا، تحلیل داده، تبدیل گفتار به متن و تولید تصویر
یافته‌اند~\cite{openai2024}.

شرکت‌های بزرگ فناوری از جمله \lr{OpenAI}، \lr{Google} (با محصول
\lr{Gemini})، \lr{Anthropic} (با محصول \lr{Claude})، \lr{Groq} و
سرویس‌های واسط مانند \lr{OpenRouter} هر کدام رابط‌های برنامه‌نویسی
منحصربه‌فرد خود را ارائه می‌دهند. گرچه بسیاری از این رابط‌ها قابلیت‌های
مشابهی دارند، اما در ساختار درخواست، احراز هویت، فرمت پاسخ و سیستم
صورت‌حساب‌گذاری تفاوت‌های قابل توجهی دارند~\cite{groq2024, openrouter2024}.

این تنوع باعث شده که سازمان‌هایی که قصد دارند از چند سرویس هوش مصنوعی
به طور هم‌زمان بهره ببرند یا از یک ارائه‌دهنده به ارائه‌دهنده‌ی دیگری مهاجرت
کنند، با چالش‌های فنی و مدیریتی قابل توجهی روبه‌رو شوند. در این شرایط،
وجود یک \term{سامانه‌ی واسط}{Gateway} که بتواند ارتباط با سرویس‌های مختلف
را به صورت استاندارد مدیریت کند، نیاز ملموسی است.

این پروژه با هدف طراحی و پیاده‌سازی چنین سامانه‌ای تعریف شده است. سامانه‌ی
پیاده‌سازی‌شده به عنوان یک لایه‌ی میانجی بین
\term{سرویس‌گیرندگان}{Clients} و ارائه‌دهندگان مختلف هوش مصنوعی عمل
می‌کند و یک رابط برنامه‌نویسی استاندارد و یکپارچه را در اختیار
توسعه‌دهندگان قرار می‌دهد. علاوه بر این، قابلیت‌هایی مانند احراز هویت،
مدیریت کلیدهای رابط برنامه‌نویسی، کیف پول، خرید \term{اشتراک}{Subscription}،
صدور فاکتور و گزارش‌های مصرف نیز در این سامانه پیش‌بینی شده است.

%----------------------------------------------------------------------
\section{بیان مسئله}
\label{sec:problem}
%----------------------------------------------------------------------

در اکوسیستم کنونی هوش مصنوعی، هر ارائه‌دهنده‌ی سرویس رابط برنامه‌نویسی
مخصوص خود را تعریف می‌کند. برای مثال، \lr{OpenAI} یک ساختار مشخص برای
ارسال پیام و دریافت پاسخ دارد، \lr{Gemini} از \lr{API} بومی \lr{Google}
با فرمت متفاوتی استفاده می‌کند، و \lr{Groq} اگرچه سازگار با فرمت \lr{OpenAI}
است، اما \lr{API Key} و \term{نقاط پایانی}{Endpoints} جداگانه‌ای دارد.
این تفاوت‌ها در چند سطح مشکل‌ساز هستند:

در سطح فنی، توسعه‌دهنده‌ای که می‌خواهد از چند سرویس هوش مصنوعی استفاده
کند ناچار است برای هر ارائه‌دهنده کد و منطق جداگانه‌ای بنویسد؛ همین موضوع
حجم کد را افزایش می‌دهد و نگهداری سامانه را دشوارتر می‌کند. در نتیجه،
مهاجرت از یک ارائه‌دهنده به دیگری نیز پرهزینه می‌شود، زیرا تغییر تنها به
نشانی سرویس یا کلید دسترسی محدود نیست و معمولاً به بازنگری در ساختار
درخواست، پاسخ و مدیریت خطا نیاز دارد.

در سطح مدیریتی و مالی نیز پراکندگی بیشتری دیده می‌شود. کلیدهای \lr{API}،
\term{محدودیت نرخ}{Rate Limit}، گزارش مصرف و صورت‌حساب هر سرویس معمولاً
به‌صورت مجزا مدیریت می‌شوند و همین امر تصویر یکپارچه‌ای از مصرف کل در
اختیار سازمان قرار نمی‌دهد. افزون بر آن، بسیاری از ارائه‌دهندگان ابزار
کافی برای مدیریت چندکاربره، کیف پول داخلی، صدور فاکتور یا خرید اشتراک
ندارند. در نبود یک لایه‌ی مرکزی، کنترل دسترسی، تعیین سقف مصرف، ممیزی
عملیات حساس و گزارش‌گیری دقیق باید برای هر سرویس به‌طور مستقل
پیاده‌سازی شود.

این مشکلات با گسترش استفاده از هوش مصنوعی در سازمان‌ها حادتر می‌شود.
راهکار پیشنهادی این پروژه، طراحی و پیاده‌سازی یک سامانه‌ی واسط است که
تمامی این چالش‌ها را از طریق یک لایه‌ی مرکزی و یکپارچه حل کند.

%----------------------------------------------------------------------
\section{انگیزه و ضرورت}
\label{sec:motivation}
%----------------------------------------------------------------------

انگیزه‌ی اصلی این پروژه از نیاز واقعی سازمان‌ها و توسعه‌دهندگانی نشأت
می‌گیرد که می‌خواهند از قابلیت‌های هوش مصنوعی بهره ببرند اما نمی‌خواهند
درگیر پیچیدگی‌های یکپارچه‌سازی چندین سرویس مجزا شوند. ضرورت چنین
سامانه‌ای را می‌توان از چند منظر بررسی کرد:

\textbf{از منظر فنی،} با داشتن یک رابط برنامه‌نویسی استاندارد و سازگار با
\lr{OpenAI} - که پرکاربردترین فرمت در این حوزه است - توسعه‌دهندگان می‌توانند
بدون تغییر در کد سمت مشتری، از ارائه‌دهندگان مختلف استفاده کنند. تنها با
تغییر فیلدهای \code{provider} و \code{model} در درخواست، سامانه
درخواست را به ارائه‌دهنده‌ی مناسب هدایت می‌کند.

\textbf{از منظر مالی،} یک سیستم متمرکز مدیریت کیف پول و صورت‌حساب،
امکان صدور فاکتور یکپارچه، ردیابی دقیق مصرف و کنترل هزینه‌ها را فراهم
می‌کند. این موضوع برای سازمان‌هایی که هزینه‌های هوش مصنوعی بخش قابل
توجهی از بودجه‌ی عملیاتی آن‌ها را تشکیل می‌دهد، از اهمیت ویژه‌ای
برخوردار است.

\textbf{از منظر امنیتی،} متمرکزسازی احراز هویت و مدیریت کلیدهای
\lr{API} در یک سامانه‌ی واحد، امکان اعمال سیاست‌های امنیتی منسجم - مانند
محدودیت نرخ درخواست، لیست سفید آدرس \lr{IP} و
\term{ممیزی}{Audit} عملیات حساس - را به صورت مرکزی فراهم می‌آورد.

\textbf{از منظر مقیاس‌پذیری،} \term{الگوی خط لوله}{Pipeline Pattern}
به‌کاررفته در این سامانه، امکان افزودن مراحل جدید به پردازش درخواست - مثلاً
\term{مدار قطعه}{Circuit Breaker} یا \term{راهبری هوشمند}{Smart Routing} -
را بدون تغییر در ساختار اصلی فراهم می‌کند.

%----------------------------------------------------------------------
\section{اهداف پروژه}
\label{sec:objectives}
%----------------------------------------------------------------------

هدف اصلی این پروژه، طراحی و پیاده‌سازی یک سامانه‌ی واسط جامع برای
یکپارچه‌سازی سرویس‌های هوش مصنوعی است. این هدف کلی به اهداف عملیاتی
زیر تجزیه می‌شود. نخست، سامانه باید مجموعه‌ای از نقاط پایانی سازگار با
استاندارد \lr{OpenAI} را برای \term{تکمیل گفتگو}{Chat Completions}،
\term{جاسازی متن}{Embeddings}، \term{تولید تصویر}{Image Generation}،
\term{تبدیل گفتار به متن}{Audio Transcription} و
\term{تبدیل متن به گفتار}{Text-to-Speech} در اختیار مشتری قرار دهد تا
اتصال سرویس‌گیرندگان با حداقل تغییر ممکن انجام شود. در کنار این قرارداد
بیرونی، پشتیبانی از چند ارائه‌دهنده از طریق \term{آداپتور}{Adapter}های
مجزا برای \lr{OpenAI}، \lr{Gemini}، \lr{Groq} و \lr{OpenRouter} نیز
لازم است تا افزودن ارائه‌دهندگان جدید به‌صورت
\term{اتصال‌پذیر}{Pluggable} ممکن بماند.

از سوی دیگر، پروژه فقط به لایه‌ی فنی محدود نیست. پیاده‌سازی جریان احراز
هویت مبتنی بر رمز یک‌بارمصرف برای ورود به داشبورد مدیریتی، مدیریت کلیدهای
\lr{API} برای درخواست‌های برنامه‌ای، و یک سیستم مالی شامل کیف پول،
ثبت دقیق مصرف توکن و منابع غیرتوکنی، صدور فاکتور و مدیریت اشتراک‌های
دوره‌ای از اهداف اصلی آن به‌شمار می‌روند. در نهایت، سامانه باید گزارش‌های
تفصیلی مصرف، دفترکل تراکنش‌های کیف پول و فهرست فاکتورها را تولید کند،
گزارش‌های ساخت‌یافته را به سرویس‌های بیرونی بفرستد و با تکیه بر آزمون‌های
واحد و ویژگی و مستندات \lr{OpenAPI} به اندازه‌ی کافی ارزیابی‌پذیر و
قابل نگهداشت باشد.

%----------------------------------------------------------------------
\section{پرسش‌ها و معیارهای ارزیابی}
\label{sec:research-questions}
%----------------------------------------------------------------------

برای آن‌که مسیر استدلال این پایان‌نامه از تعریف مسئله تا ارزیابی نهایی
شفاف باشد، سه پرسش اصلی زیر به‌عنوان محور ارزیابی در نظر گرفته شده‌اند:

\begin{enumerate}

  \item \textbf{پرسش اول:} آیا می‌توان یک رابط بیرونی سازگار با \lr{OpenAI}
    ارائه کرد که ناهمگنی چند ارائه‌دهنده‌ی متفاوت - به‌ویژه ارائه‌دهنده‌های
    سازگار و ناسازگار با این قرارداد - را در لایه‌ی درونی پنهان کند؟

  \item \textbf{پرسش دوم:} آیا می‌توان در کنار یکپارچه‌سازی فنی، قابلیت‌های
    مدیریتی و مالی مورد نیاز یک سازمان - مانند احراز هویت، مدیریت کلید،
    کیف پول، اشتراک، فاکتور و گزارش مصرف - را نیز در همان درگاه تجمیع کرد؟

  \item \textbf{پرسش سوم:} آیا پیاده‌سازی انجام‌شده از نظر درستی جریان‌های
    بحرانی، قابلیت توسعه و انسجام معماری، شواهد کافی برای یک نمونه‌ی
    کارکردی و قابل اتکا ارائه می‌دهد؟

\end{enumerate}

برای پاسخ به این پرسش‌ها، چهار معیار ارزیابی در این پایان‌نامه به‌کار
گرفته شده است. معیار نخست، کامل بودن قرارداد بیرونی است؛ یعنی این‌که آیا
نقاط پایانی اصلی هوش مصنوعی و مسیرهای مدیریتی مورد نیاز در سامانه
پیاده‌سازی شده‌اند یا نه. معیار دوم، درستی رفتار در سناریوهای بحرانی
است و مسیرهایی مانند احراز هویت، کنترل موجودی، صورت‌حساب، گزارش‌گیری و
نرمال‌سازی پاسخ ارائه‌دهندگان را با تکیه بر آزمون‌های خودکار می‌سنجد.
معیار سوم به انسجام معماری و قابلیت توسعه مربوط می‌شود و بررسی می‌کند که
آیا تفکیک مسئولیت‌ها، الگوی خط لوله و آداپتورها توسعه‌ی بعدی سامانه را
تسهیل می‌کنند. معیار چهارم نیز مرز اعتبار نتایج را روشن می‌کند؛ به این
معنا که این پایان‌نامه درستی نرم‌افزاری و کفایت معماری را می‌سنجد و
مدعی بنچ‌مارک میدانی گسترده یا یکپارچه‌سازی کامل با همه‌ی سرویس‌های
بیرونی نیست.

%----------------------------------------------------------------------
\section{دستاوردهای پروژه}
\label{sec:contributions}
%----------------------------------------------------------------------

پروژه‌ی حاضر چند دستاورد مشخص را به‌عنوان خروجی اصلی ارائه می‌دهد. مهم‌ترین
آن‌ها یک سامانه‌ی واسط کارکردی در سمت سرور است که درخواست‌های هوش مصنوعی
را از سرویس‌گیرندگان دریافت و به چهار ارائه‌دهنده‌ی \lr{OpenAI}،
\lr{Gemini}، \lr{Groq} و \lr{OpenRouter} هدایت می‌کند؛ شواهد ارزیابی این
ادعا در فصل پنجم آمده است. در کنار آن، یک رابط برنامه‌نویسی سازگار با
\lr{OpenAI} پیاده‌سازی شده است تا سرویس‌گیرندگان موجود با حداقل تغییر به
سامانه متصل شوند.

از نظر درونی، معماری سامانه بر چهارده مرحله‌ی پردازشی مستقل در قالب خط
لوله‌ی \lr{Laravel} استوار است و همین موضوع توسعه و آزمون آن را ساده‌تر
می‌کند. همچنین سیستم صورت‌حساب با استفاده از
\term{توکن‌ساز واقعی}{Real BPE Tokenizer} برای سنجش مصرف، پشتیبانی از
صورت‌حساب‌گذاری منابع غیرتوکنی مانند تصویر و صوت، و مدیریت کیف پول با
تراکنش‌های ایمن تکمیل شده است. \lr{Playground} مبتنی بر \lr{Vue.js} برای
آزمایش تعاملی نقاط پایانی و مجموعه‌ای از آزمون‌های واحد و ویژگی نیز این
خروجی را از حالت صرفاً نمایشی خارج کرده و به یک کدپایه‌ی قابل ارزیابی
تبدیل کرده‌اند.

%----------------------------------------------------------------------
\section{تغییرات نسبت به پروپوزال}
\label{sec:changes}
%----------------------------------------------------------------------

در \term{پروپوزال}{Proposal} اولیه‌ی این پروژه، بسته‌ی فناوری پیش‌بینی‌شده
شامل زبان \lr{Python}، چارچوب‌های \lr{FastAPI} یا \lr{Django}، و پایگاه
داده‌ی \lr{PostgreSQL} بود. در مرحله‌ی اجرا، این بسته‌ی فناوری تغییر یافت و
پیاده‌سازی نهایی با فناوری‌های \lr{Laravel} بر پایه‌ی \lr{PHP} در سمت
سرور، \lr{SQL Server} برای ذخیره‌سازی داده، \lr{Redis} برای
\term{موتور کش}{Cache Engine} و مدیریت صف‌های پس‌زمینه، و \lr{Vue.js}
برای پیاده‌سازی \term{میدان بازی}{Playground} انجام شد.

اهداف اصلی پروژه - شامل یکپارچه‌سازی چند ارائه‌دهنده‌ی هوش مصنوعی، سیستم
احراز هویت، مدیریت کلیدهای \lr{API}، کیف پول، اشتراک و گزارش‌دهی - بدون
تغییر باقی ماندند و کاملاً پیاده‌سازی شدند.

%----------------------------------------------------------------------
\section{ساختار پایان‌نامه}
\label{sec:structure}
%----------------------------------------------------------------------

مابقی این پایان‌نامه به این صورت سازماندهی شده است: در فصل دوم، مفاهیم
پایه‌ی سامانه‌های واسط و رابط‌های برنامه‌نویسی هوش مصنوعی بررسی می‌شوند و
سپس نمونه‌های مشابهی مانند \lr{Groq}، \lr{OpenRouter}،
\lr{Anyscale Endpoints} و \lr{OpenAI} تحلیل و مقایسه می‌شوند. فصل سوم به
طراحی سامانه اختصاص دارد و معماری کلی، مدل داده، طراحی \lr{API}، الگوی
خط لوله، استراتژی کش و سیستم صورت‌حساب را شرح می‌دهد. در فصل چهارم،
جزئیات پیاده‌سازی ماژول‌های اصلی از جمله خط لوله‌ی سامانه‌ی واسط،
آداپتورهای ارائه‌دهندگان، احراز هویت، مدیریت کلیدهای \lr{API} و
صورت‌حساب ارائه می‌شود. فصل پنجم نتایج آزمون‌های واحد و ویژگی، ارزیابی
عملکرد و مقایسه با اهداف اولیه را جمع‌بندی می‌کند و فصل ششم به نتیجه‌گیری
و پیشنهاد کارهای آتی می‌پردازد.
